\documentclass[10pt,a4paper]{amsart}
\usepackage[margin=19mm]{geometry}
\usepackage{amsmath,amssymb,mathtools}
\usepackage[scr=rsfs,cal=euler]{mathalfa}
\usepackage{xfrac}
\usepackage[unicode,hidelinks,bookmarks=false]{hyperref}
\newcommand{\ie}{i.e.\ }
\newcommand{\cf}{cf.\ }
\newcommand{\resp}{respectively\ }
\newcommand{\Subsection}{\subsection}
\providecommand{\nobreakdash}{\nobreak\hbox{-}}
\setlength{\emergencystretch}{2em}
\multlinegap0pt
\usepackage{amssymb}
\usepackage{dsfont}
\newtheorem{lemma}{Lemma}
\DeclareMathOperator{\diam}{diam}
\title{Piecewise contracting maps on the interval: Hausdorff dimension, entropy and attractors}
\author{A.E. Calder\'on, E. Villar-Sep\'ulveda}
\date{Source: arXiv:2309.09154v1 (CC BY 4.0)}
\begin{document}
\maketitle

\noindent\emph{Source and licence.} Piecewise contracting maps on the interval: Hausdorff dimension, entropy and attractors. Version arXiv:2309.09154v1 (CC BY 4.0), distributed by arXiv under the Creative Commons Attribution 4.0 International licence, http://creativecommons.org/licenses/by/4.0/, as recorded in the arXiv licence metadata for this exact version and shown on the abstract page https://arxiv.org/abs/2309.09154v1. Full licence notice: \texttt{licenses/arxiv-2309.09154-CC-BY-4.0.txt}.\par\smallskip

\noindent\textit{Editorial scope.} Counted unit: the article's lemma WELL (source0.tex lines 293--295). The article's complete original proof of it is delivered with it (source0.tex lines 296--306, one \texttt{proof} environment of 11 lines). No mathematics is written, repaired or re-derived: the statement and the proof are the article's own lines, in the article's own notation, and no retained line is substituted or re-flowed.  Retained uncounted as the source-local context the proof needs: lines 163--165; lines 179--193.  Nothing was omitted from any retained span, so the package carries no omission record.  No retained line contains a citation, so the wrapper carries no bibliography and no bibliography entry.  No retained line contains a figure environment, an image-inclusion command or drawing code, so no image is delivered and the package declares no asset dependency.  Editorial formatting only, in addition: an AMS \texttt{amsart} preamble that restates the article's own declarations and macros used by the retained lines, namely the environments \texttt{lemma} on separate counters, together with the standard packages those lines need.  The retained environments print in this document as Lemma 1 in order of appearance, whereas the article prints the same items with the numbering of its own full document.  The complete native source of the article as distributed in the arXiv e-print for this exact version is bundled unchanged as \texttt{sources/147140/source0.tex}.
\begin{equation}
    \Lambda=\bigcap_{n\geq1}\Lambda_n\,,\quad \text{ where }\;\, \Lambda_n=\bigcup_{A\in\mathcal A_n}A\quad\forall\,n\geq1.\label{ATTRACTOR.ATOMS}
\end{equation}

\begin{lemma}\label{ATOMS.LEMMA}
    Each one of the following statements holds:
    \begin{enumerate}
        \item For all $n\geq1$,
        \[
            \max_{A\in\mathcal A_{n+1}}\diam(A)\leq \lambda \max_{A\in\mathcal A_{n}}\diam(A)\,;
        \]
        \item For all $\epsilon>0$, there exists $n_0\geq1$ such that $\diam(A)<\epsilon$ for every $A\in\mathcal A_{n}$ with $n\geq n_0\,$;
        \item For all $x\in\Lambda$, there exists a decreasing sequence $\{B_k\}_{k\geq1}$ of atoms (i.e. $B_k\supset B_{k+1}$ for all $k\geq 1$) such that $x\in B_k\in\mathcal A_k$ for all $k\geq1$ and
        \begin{equation}
            \bigcap_{k\geq1}B_k=\{x\}\,;\label{ONEPOINT}
        \end{equation}
        \item If $f$ satisfies the separation property, the sequence  $\{B_n\}_{n\geq1}$ defined in item 3 is unique for each $x\in\Lambda$.
    \end{enumerate}
\end{lemma}

\begin{lemma}\label{WELL}
    Suppose that $f$ is piecewise monotonic and $D\subset\widetilde X$. Then, for every $n\geq1$ and $\epsilon>0$, $r_n(\epsilon,\Lambda\cap\widetilde X)$ is a finite number.
\end{lemma}

\begin{proof}
Let $n\geq1$ and $\epsilon>0$. From item 2 of Lemma \ref{ATOMS.LEMMA}, there exists $n_0\geq n$ such that
    \begin{equation}
        \diam A<\epsilon\qquad \forall\,A\in\mathcal A_{n_0}. \label{ATOM.EPSILON}
    \end{equation}
    Next, denote by $\mathcal C(\mathcal A_{n_0},n)$ the collection of all connected components of
    \[
        A\setminus\big(\Delta\cup f^{-1}(\Delta)\cup\ldots\cup f^{-n+1}(\Delta)\big)\quad\text{with $A \in \mathcal A_{n_0}$.}
    \]
{Since $f$ is piecewise monotonic, the collection $\mathcal C(\mathcal A_{n_0},n)$ is finite. Furthermore, from \eqref{ATTRACTOR.ATOMS} we can consider the subcollection $\mathcal C^*(\mathcal A_{n_0},n)\subset\mathcal C(\mathcal A_{n_0},n)$ of connected components that intersect $\Lambda\cap\widetilde X$, whose union covers said set. Thus, for every $B\in \mathcal C^*(\mathcal A_{n_0},n)$ we can take $x_B\in B\cap\Lambda\cap\widetilde X\neq\emptyset$. It is clear that $E_{n,\epsilon}:=\{x_B\ :\ B\in\mathcal C^*(\mathcal A_{n_0},n)\}\subset\Lambda\cap\widetilde X$ is a finite $(n,\epsilon)$-spanning set. This implies that $r_n(\epsilon,\Lambda\cap\widetilde X)<\infty$, concluding the proof.}
\end{proof}

\end{document}
